\clearpage

\section{Conclusion and Outlook}
  \label{sec:conclusion-and-outlook}

  Cosmochrony reduces the fundamental assumptions of physics to a single
  structural origin: the non-injective projection of a static pre-geometric
  relational substrate~$\chi$ onto an observable effective description.
  From this substrate, time, spacetime geometry, and a wide spectrum of
  physical phenomena emerge as effective descriptions obtained through
  coarse-grained, generally non-injective projection of admissible relational configurations.

  A central structural ingredient of the framework is the spectral
  organization of the admissible sector of the substrate.
  The projective structure induces a spectral organization on the space of
  $\chi$ configurations, and physical states correspond to projectable
  modes selected by admissibility constraints.
  This spectral filtering provides the bridge between the relational
  substrate and the effective hierarchy of physical excitations,
  and underlies the emergence of mass scales through admissible spectral localisation.

  Within this setting, bounded-response regimes arise from saturation of
  the admissible projective flux.
  The resulting effective dynamics take a Born--Infeld-like form without
  introducing additional fundamental fields.
  The speed limit~$c_\chi$ and Planck's constant~$h$ appear as
  complementary bounds on projectability, respectively limiting the maximal
  propagation rate and the minimal resolvable granularity.
  General Relativity emerges as a thermodynamic and geometric limit of the underlying
  relational dynamics in regimes where local injectivity is approximately recovered.

  The framework also provides a structural interpretation of the origin of
  matter and interactions.
  Gauge interactions arise as aspects of projection dynamics associated with
  the topology of the projection fiber, while mass and particle stability are
  associated with localized spectral configurations and topological constraints
  of the admissible sector.
  Electric charge is identified with a $\pi_1$ winding invariant of the fiber,
  while the absence of magnetic monopoles follows from the admissibility-induced
  triviality of $\pi_2$ on the canonical admissible base.
  Quantum indeterminacy and entanglement appear as generic consequences of
  non-injective projection rather than as signatures of fundamental
  nonlocal dynamics.

  At cosmological scales, expansion and the arrow of time follow from the
  global growth of cumulative projected capacity~$I(n)$~\cite{Beau2026pto}.
  At galactic scales, saturation of the admissible projective capacity produces
  flat rotation curves without requiring dark matter halos.
  In strong-gravity regimes, black hole evaporation can be interpreted as
  a reprojection process ensuring information preservation within the
  relational description, without invoking fundamental information loss.

  The space of admissible projected configurations evolves as projective
  ordering advances: the particle content of the Universe is therefore a
  historically conditioned outcome rather than a fixed input.

  \subsection*{Conceptual Shift and Outlook}

    Cosmochrony represents a shift from a ``matter-on-spacetime'' paradigm
    to a ``projection-of-substrate'' ontology, in which spacetime,
    quantum behaviour, and interactions emerge from a single underlying
    structural fact: the non-injective projection~$\Pi$ of a static
    relational substrate~$\chi$.

    Immediate falsifiable research directions include:
    cosmological signatures (low-$\ell$ CMB anomalies induced by finite projective
    capacity), galactic phenomenology (systematic correlation between apparent
    dark matter effects and local projective capacity gradients), and
    fundamental-scale effects (environment-dependent variations in decay rates
    or effective couplings in regimes approaching admissibility saturation).

    On the theoretical side, a key open problem is the derivation of the full
    Standard Model structure as a hierarchy of admissible spectral excitations
    and topological classes within the relational projective admissibility
    structure, including the quantitative emergence of chirality selection and
    the $V{-}A$ structure of weak interactions.
